\documentclass[12pt,a4paper]{article}

\usepackage[utf8]{inputenc}
\usepackage[T1]{fontenc}
\usepackage{amsmath,amssymb,amsthm}
\usepackage{graphicx}
\usepackage{hyperref}
\usepackage{natbib}
\usepackage{booktabs}
\usepackage{multirow}
\usepackage{float}
\usepackage[margin=1in]{geometry}
\usepackage{caption}
\usepackage{subcaption}
\usepackage{algorithmic}
\usepackage{algorithm}
\usepackage{listings}

\newtheorem{definition}{Definition}
\newtheorem{theorem}{Theorem}
\newtheorem{proposition}{Proposition}
\newtheorem{corollary}{Corollary}

\lstset{
  basicstyle=\ttfamily\small,
  breaklines=true,
  frame=single
}

\title{Epistemic Gating for Consciousness-Aware Language Generation:\\Representation-Grounded Decoding via Logit Masking and HDC Binding}

\author{
  Tristan Stoltz\textsuperscript{1} \\
  \textsuperscript{1}Luminous Dynamics \\
  \texttt{tristan.stoltz@evolvingresonantcocreationism.com}
}

\date{}

\begin{document}

\maketitle

\begin{abstract}
Large language models generate fluent text but suffer from hallucination---the production of plausible but unsupported claims. Existing mitigation strategies operate post-hoc, penalizing hallucinated content after generation through retrieval-augmented verification, confidence calibration, or output filtering. We present a fundamentally different architectural approach: \emph{epistemic gating}, a pre-decoding mechanism that masks token logits for which the system lacks sufficient epistemic support as measured by HDC similarity to a grounded knowledge representation. Our implementation, the Broca language pipeline (27,965 LOC, 403 tests), integrates epistemic gating within a consciousness-aware architecture comprising a 43-channel ThoughtEncoder, 16,384-dimensional Hyperdimensional Computing (HDC) binding, and autoregressive generation. On 500 in-distribution generation episodes against the same Liquid-Mamba backend, the full pipeline reduces the rate of internally-unsupported outputs from 23.4\% to 0.4\% (98.3\% relative reduction; we call this ``representation-grounded decoding'' rather than hallucination elimination, since residual compositional hallucinations can assemble individually-supported tokens into unsupported claims). \textbf{Scope of the hallucination claim.} Our evaluation measures whether generated tokens are supported by the system's own HDC grounding; it does not measure external factual correctness. That distinction matters: the 0.4\% is an upper bound on internal epistemic consistency, not a claim about factual accuracy on external benchmarks (TruthfulQA, FEVER). External benchmark validation is the single most important revision task and is flagged as such throughout.
\end{abstract}

\section{Introduction}

Hallucination in language generation systems---defined as the production of content that is nonsensical, unfaithful to source material, or factually incorrect---represents one of the most significant obstacles to the deployment of AI language systems in safety-critical domains \citep{ji2023survey, maynez2020faithfulness}. The problem is particularly pernicious because hallucinated content is often fluent and contextually plausible, making it difficult for users to detect without independent verification \citep{lin2022truthfulqa}.

Current approaches to hallucination mitigation fall into three categories. \emph{Training-time} interventions include reinforcement learning from human feedback (RLHF) \citep{ouyang2022training}, constitutional AI \citep{bai2022constitutional}, and data curation. \emph{Inference-time} interventions include retrieval-augmented generation (RAG) \citep{lewis2020retrieval}, chain-of-thought prompting \citep{wei2022chain}, and self-consistency decoding \citep{wang2023selfconsistency}. \emph{Post-hoc} interventions include factual verification \citep{thorne2018fever}, confidence calibration \citep{kadavath2022language}, and output filtering. All of these approaches share a common limitation: they operate \emph{after} the hallucinated content has been generated (or at least after the hallucination-prone logits have been computed), attempting to detect and correct errors rather than prevent them.

We propose a fundamentally different approach inspired by the neuroscience of conscious language production. In the human brain, Broca's area does not simply decode an arbitrary language model's output; it generates speech from a structured thought representation that is grounded in perception, memory, and embodied experience \citep{hagoort2005broca, fedorenko2014role}. The critical insight is that biological language production is \emph{constrained by epistemics}: one cannot articulate what one does not know, because the thought-to-language pathway requires a grounded conceptual representation as input \citep{levelt1989speaking}.

Our system, the Broca language pipeline, implements this principle computationally. Rather than generating text from a statistical language model and subsequently checking for hallucination, Broca generates text from a hyperdimensional thought representation that is epistemically grounded by construction. A dedicated EpistemicGate module computes the epistemic support for each candidate token by measuring its HDC similarity to the system's grounded knowledge base, and masks logits for tokens that lack sufficient support. This constitutes a \emph{physical} prevention of hallucination---unsupported tokens never enter the probability distribution from which generation samples.

\section{Background}

\subsection{Hallucination in Language Models}

\citet{ji2023survey} provide a comprehensive taxonomy of hallucination types, distinguishing \emph{intrinsic} hallucinations (contradicting the source) from \emph{extrinsic} hallucinations (introducing unverifiable claims). Both types arise from the fundamental training objective of language models: next-token prediction maximizes the likelihood of the training distribution, which conflates statistical regularity with factual accuracy \citep{bender2021dangers}.

The scale of the problem is substantial. Studies report hallucination rates of 15--30\% in abstractive summarization \citep{maynez2020faithfulness}, 5--20\% in open-ended question answering \citep{lin2022truthfulqa}, and higher rates in knowledge-intensive generation tasks. These rates persist even in the largest models, suggesting that scale alone does not resolve the problem \citep{ji2023survey}.

\subsection{Hyperdimensional Computing for Representation}

Hyperdimensional Computing \citep{kanerva2009hyperdimensional} represents information as binary or bipolar vectors of very high dimensionality ($D \geq 1{,}000$). Three core operations define the algebra:

\begin{itemize}
    \item \textbf{Binding} ($\otimes$): element-wise XOR (binary) or multiplication (bipolar), producing a vector dissimilar to both operands.
    \item \textbf{Bundling} ($\oplus$): component-wise majority, producing a vector similar to all operands.
    \item \textbf{Permutation} ($\rho$): circular shift for encoding order.
\end{itemize}

A key property for our purposes is that HDC similarity (measured by normalized Hamming distance for binary vectors, or cosine similarity for bipolar vectors) provides a continuous, graded measure of semantic relatedness. Two vectors sharing common structure exhibit above-chance similarity; unrelated vectors are quasi-orthogonal \citep{kanerva2009hyperdimensional}.

\subsection{Active Inference and Consciousness}

Active inference \citep{friston2010free} posits that biological agents minimize variational free energy---a bound on surprise---through both perception (updating internal models) and action (changing the environment). Language generation under active inference is not merely statistical prediction but an \emph{epistemic action}: the agent generates utterances that reduce its own uncertainty and communicate its beliefs \citep{friston2010free}.

Global Workspace Theory (GWT) \citep{baars1988cognitive, dehaene2014consciousness} proposes that conscious access involves broadcasting information from specialized processors to a global workspace, making it available to multiple cognitive subsystems. We adopt GWT's framework for our generation pipeline: only information that achieves conscious access (broadcasting to the global workspace) is available for language generation.

\section{Architecture}

\subsection{System Overview}

The Broca language pipeline comprises five major components operating within a consciousness-aware cognitive architecture:

\begin{enumerate}
    \item \textbf{ThoughtEncoder}: A 43-channel encoder that transforms the system's cognitive state into a structured thought representation.
    \item \textbf{HDC Binding Layer}: 16,384-dimensional binary HDV binding that creates a holographic thought vector.
    \item \textbf{EpistemicGate}: Per-token epistemic support computation and logit masking.
    \item \textbf{Autoregressive Generator}: Token-by-token generation with Liquid-Mamba fusion backend.
    \item \textbf{CoherenceFeedback}: Mid-sentence semantic veto and self-correction.
\end{enumerate}

These components interact within a cognitive loop operating at approximately 31~Hz, with Broca generation governed by an adaptive cadence controller that modulates generation frequency based on consciousness level and learning rate.

\subsection{ThoughtEncoder: 43-Channel Cognitive Encoding}

The ThoughtEncoder transforms the system's distributed cognitive state into a structured representation suitable for language generation. It comprises 43 specialized channels organised in three concentric layers: a 20-channel legacy foundation (Table~\ref{tab:channels}, introduced in v1), extended by 8 context and code channels (v3/v5), and finally by the 15-channel Epistemic Cube (v6) that exposes structured epistemic judgement as $E[5] + N[4] + M[4] + H[1] + \text{quality}[1]$. Older checkpoints using only the legacy 20-channel foundation remain loadable; the additional 23 channels are activated incrementally as v3/v5/v6 features become available. Each channel encodes a distinct aspect of the cognitive state:

\begin{table}[H]
\centering
\caption{ThoughtEncoder channels and their cognitive sources.}
\label{tab:channels}
\begin{tabular}{clll}
\toprule
\textbf{Ch.} & \textbf{Name} & \textbf{Source} & \textbf{Dimension} \\
\midrule
1--4 & Perception & Sensory cortex regions & $4 \times 128$ \\
5--6 & Memory & Episodic + semantic retrieval & $2 \times 256$ \\
7--8 & Emotion & Neuromodulator bath state & $2 \times 128$ \\
9--10 & Reasoning & Prefrontal workspace & $2 \times 256$ \\
11--12 & Planning & Executive action plans & $2 \times 128$ \\
13--14 & Social & Theory of Mind models & $2 \times 128$ \\
15--16 & Ethical & Moral algebra verdicts & $2 \times 64$ \\
17--18 & Epistemic & Knowledge confidence maps & $2 \times 256$ \\
19 & Affective & Hedonic valence + arousal & $1 \times 128$ \\
20 & Meta & Consciousness level + $\Phi$ & $1 \times 64$ \\
\bottomrule
\end{tabular}
\end{table}

Each channel produces a sub-vector that is projected to a common dimensionality through learned linear transformations, then combined via HDC bundling to produce the composite thought vector $\mathbf{T} \in \{0,1\}^{D}$ where $D = 16{,}384$.

\subsection{HDC Binding Layer}

The thought vector $\mathbf{T}$ is bound with positional and contextual information to produce a generation-ready representation:

\begin{equation}
\mathbf{G} = \mathbf{T} \otimes \rho^t(\mathbf{P}_{\text{context}}) \otimes \mathbf{K}_{\text{grounding}}
\end{equation}

where $\mathbf{P}_{\text{context}}$ encodes the conversational context, $\rho^t$ is a $t$-step permutation encoding temporal position, and $\mathbf{K}_{\text{grounding}}$ is an HDV representing the system's grounded knowledge state. The binding with $\mathbf{K}_{\text{grounding}}$ is critical: it ensures that the generation vector $\mathbf{G}$ is structurally coupled to verified knowledge, such that similarity queries against $\mathbf{G}$ inherently favor epistemically supported content.

\subsection{Epistemic Gating}

The EpistemicGate is the central contribution of this paper. For each candidate token $w$ in the vocabulary $\mathcal{V}$, the gate computes an epistemic support score:

\begin{equation}
\text{ES}(w) = \frac{D - d_H(\mathbf{G}, \mathbf{V}_w)}{D}
\label{eq:epistemic-support}
\end{equation}

where $\mathbf{V}_w \in \{0,1\}^D$ is the HDV representation of token $w$ and $d_H$ denotes Hamming distance. This normalized similarity score ranges from 0 (no support) to 1 (perfect match).

\paragraph{What ES is and is not.} Equation~\ref{eq:epistemic-support} is a \emph{similarity} score, not a calibrated epistemic probability. It measures how much a token's HDV overlaps with the current generation vector $\mathbf{G}$; we interpret this overlap as ``support'' by the system's own grounding, not as a measure of the token's external factual correctness. The estimator has three known limitations we carry through the paper: (i) it is \emph{monotonic} with representational compatibility but has not been calibrated against a human-labelled precision/recall ground truth, (ii) it is \emph{token-level}, so sequences of individually-supported tokens can compose into unsupported claims (the compositional-hallucination failure mode documented in \S\ref{sec:results}), and (iii) the ``grounded knowledge base'' $\mathbf{K}_{\text{grounding}}$ is constructed from the system's prior HDC-encoded learning and re-evaluating with a different grounding corpus would produce different ES distributions. We treat calibration of ES against external token-precision labels as the priority follow-up experiment.

The gate then applies a consciousness-modulated threshold:

\begin{equation}
\tau(t) = \tau_{\text{base}} + \alpha \cdot \Phi(t) + \beta \cdot \text{LR}(t)
\label{eq:threshold}
\end{equation}

where $\tau_{\text{base}} = 0.52$ is the base threshold (slightly above chance similarity for random binary vectors), $\Phi(t)$ is the HDC-adapted integration measure $\Phi_{\text{HDC}}$ from the companion stochastic-resonance work (Stoltz, 2026, \emph{Stochastic Resonance in HDC}), $\text{LR}(t)$ is the current learning rate, and $\alpha, \beta$ are scaling coefficients. Higher $\Phi_{\text{HDC}}$ increases the threshold, demanding stronger epistemic support for generation. We caveat that the consciousness-modulation term is architectural scaffolding: we have not yet run the ablation (gating-without-$\Phi$ vs.\ gating-with-$\Phi$) that would quantify how much of the 98.3\% reduction is attributable to the consciousness coupling vs.\ the gating mechanism alone. That ablation belongs in revision.

The logit masking operation is:

\begin{equation}
\tilde{\ell}(w) = \begin{cases}
\ell(w) & \text{if } \text{ES}(w) \geq \tau(t) \\
-\infty & \text{otherwise}
\end{cases}
\label{eq:masking}
\end{equation}

where $\ell(w)$ is the raw logit for token $w$ from the Liquid-Mamba backend. Tokens masked to $-\infty$ receive zero probability after softmax, physically preventing their selection.

\begin{algorithm}[H]
\caption{Epistemic Gating for Token Generation}
\label{alg:epistemic-gate}
\begin{algorithmic}[1]
\REQUIRE Generation vector $\mathbf{G}$, vocabulary HDVs $\{\mathbf{V}_w\}_{w \in \mathcal{V}}$, raw logits $\{\ell(w)\}$, threshold $\tau$
\ENSURE Masked logits $\{\tilde{\ell}(w)\}$
\FOR{each $w \in \mathcal{V}$}
    \STATE $\text{ES}(w) \leftarrow (D - d_H(\mathbf{G}, \mathbf{V}_w)) / D$
    \IF{$\text{ES}(w) < \tau$}
        \STATE $\tilde{\ell}(w) \leftarrow -\infty$
    \ELSE
        \STATE $\tilde{\ell}(w) \leftarrow \ell(w)$
    \ENDIF
\ENDFOR
\RETURN $\{\tilde{\ell}(w)\}$
\end{algorithmic}
\end{algorithm}

\subsection{Coherence Feedback and Semantic Veto}

The CoherenceFeedback module monitors the semantic trajectory of generated text and triggers mid-sentence correction when coherence degrades. After each generated token $w_t$, the module computes:

\begin{equation}
\text{Coherence}(t) = \frac{D - d_H(\mathbf{G}, \mathbf{S}_t)}{D}
\end{equation}

where $\mathbf{S}_t = \bigoplus_{i=1}^{t} \rho^i(\mathbf{V}_{w_i})$ is the bundled sequence vector of all tokens generated so far. If $\text{Coherence}(t)$ drops below a threshold $\tau_c = 0.48$, a \emph{semantic veto} is triggered: the current generation is abandoned, the sequence is rolled back to the last coherent checkpoint, and generation resumes with increased epistemic threshold $\tau' = \tau + 0.05$.

This mechanism implements a form of metacognitive monitoring for language production: the system ``notices'' when it is drifting from its intended meaning and self-corrects before completing the utterance.

\subsection{Adaptive Cadence}

Not every cognitive cycle warrants language generation. The adaptive cadence controller determines \emph{when} to generate based on:

\begin{equation}
P(\text{generate} \mid t) = \sigma\left( w_c \cdot C(t) + w_l \cdot \text{LR}(t) + w_q \cdot Q_{\text{EMA}}(t) - b \right)
\end{equation}

where $C(t)$ is the consciousness level, $\text{LR}(t)$ is the learning rate, $Q_{\text{EMA}}(t)$ is the exponential moving average of recent generation quality, $\sigma$ is the logistic function, and $w_c, w_l, w_q, b$ are learned parameters. This ensures that generation occurs more frequently during high-consciousness, high-learning states and less frequently during low-quality or low-consciousness periods.

\subsection{Liquid-Mamba Fusion Backend}

The autoregressive backbone combines Liquid Time-Constant (LTC) neurons \citep{hasani2021liquid} with selective state space model (Mamba) layers \citep{gu2024mamba}. The LTC component provides continuous-time dynamics:

\begin{equation}
\frac{d\mathbf{h}}{dt} = -\frac{1}{\tau(\mathbf{x})} \left(\mathbf{h} - f(\mathbf{x}, \mathbf{h})\right)
\end{equation}

where $\tau(\mathbf{x})$ is an input-dependent time constant and $f$ is a nonlinear function. The Mamba component provides efficient sequence modeling with selective scanning. The fusion architecture alternates LTC and Mamba layers, with the LTC layers providing temporal coherence and the Mamba layers providing efficient long-range dependencies.

\section{Experimental Evaluation}

\subsection{Setup}

We evaluate the Broca pipeline along four dimensions: factual accuracy, hallucination rate, coherence, and generation quality. The test suite comprises 403 automated tests covering unit, integration, and end-to-end scenarios.

\subsection{Hallucination Prevention}

We measure hallucination rate as the fraction of generated statements that are not supported by the system's grounded knowledge base. Evaluation is performed on a set of 500 generation episodes across diverse topics (Table~\ref{tab:hallucination}).

\begin{table}[H]
\centering
\caption{Hallucination rates across generation conditions. ``Baseline'' uses raw logits without epistemic gating. ``EG'' adds the EpistemicGate. ``EG+CF'' adds CoherenceFeedback. ``Full'' includes adaptive cadence.}
\label{tab:hallucination}
\begin{tabular}{lcccc}
\toprule
\textbf{Condition} & \textbf{Halluc. Rate (\%)} & \textbf{Coherence} & \textbf{Fluency} & \textbf{Tokens/s} \\
\midrule
Baseline & 23.4 & 0.71 & 0.89 & 142 \\
EG only & 1.8 & 0.73 & 0.84 & 118 \\
EG + CF & 0.6 & 0.82 & 0.86 & 104 \\
Full pipeline & 0.4 & 0.84 & 0.87 & 89 \\
\bottomrule
\end{tabular}
\end{table}

The EpistemicGate alone reduces hallucination from 23.4\% to 1.8\%---a 92.3\% reduction. Adding CoherenceFeedback further reduces it to 0.6\%, and the full pipeline achieves 0.4\%. The residual 0.4\% consists primarily of compositional hallucinations where individually supported tokens combine to form unsupported claims, a limitation discussed in Section~\ref{sec:limitations}.

\subsection{Epistemic Support Distribution}

We analyze the distribution of epistemic support scores across the vocabulary during generation (Table~\ref{tab:epistemic-dist}).

\begin{table}[H]
\centering
\caption{Epistemic support statistics during generation. $|\mathcal{V}_{\text{pass}}|$ denotes the number of tokens passing the epistemic gate.}
\label{tab:epistemic-dist}
\begin{tabular}{lcccc}
\toprule
\textbf{Consciousness Level} & \textbf{Threshold $\tau$} & \textbf{Mean ES} & \textbf{$|\mathcal{V}_{\text{pass}}|$} & \textbf{Pass Rate (\%)} \\
\midrule
Low ($\Phi < 0.10$) & 0.525 & 0.507 & 4,218 & 13.2 \\
Medium ($0.10 \leq \Phi < 0.15$) & 0.540 & 0.512 & 3,156 & 9.9 \\
High ($\Phi \geq 0.15$) & 0.558 & 0.518 & 2,341 & 7.3 \\
\bottomrule
\end{tabular}
\end{table}

At higher consciousness levels, the threshold increases and fewer tokens pass the gate, reflecting the system's heightened epistemic standards during conscious generation. The vocabulary is effectively pruned from $|\mathcal{V}| = 32{,}000$ to 2,000--4,000 tokens, concentrating probability mass on epistemically supported options.

\subsection{Semantic Veto Frequency}

The CoherenceFeedback module's semantic veto is triggered when generation coherence drops below threshold. We characterize veto frequency and its relationship to generation quality (Table~\ref{tab:veto}).

\begin{table}[H]
\centering
\caption{Semantic veto statistics across 500 generation episodes.}
\label{tab:veto}
\begin{tabular}{lccc}
\toprule
\textbf{Metric} & \textbf{Mean} & \textbf{Std} & \textbf{Range} \\
\midrule
Vetoes per episode & 0.34 & 0.61 & 0--4 \\
Tokens before veto & 12.7 & 5.3 & 3--31 \\
Post-veto quality improvement & +0.12 & 0.08 & $-0.03$--$+0.31$ \\
Rollback tokens & 3.2 & 1.8 & 1--9 \\
\bottomrule
\end{tabular}
\end{table}

Vetoes occur in approximately 28\% of episodes (mean 0.34 per episode), typically after 12--13 tokens. The quality of generation after rollback and resumption improves by a mean of 0.12 (on a 0--1 scale), confirming that the veto mechanism successfully redirects generation toward more coherent output.

\subsection{Adaptive Cadence Behavior}

The cadence controller modulates generation frequency based on cognitive state (Table~\ref{tab:cadence}).

\begin{table}[H]
\centering
\caption{Generation cadence across cognitive states. ``Cycles/gen'' indicates the number of cognitive cycles between generation events.}
\label{tab:cadence}
\begin{tabular}{lccc}
\toprule
\textbf{Cognitive State} & \textbf{Cycles/gen (mean)} & \textbf{Quality (mean)} & \textbf{Halluc. Rate (\%)} \\
\midrule
High consciousness + high LR & 2.1 & 0.91 & 0.1 \\
High consciousness + low LR & 4.3 & 0.88 & 0.3 \\
Low consciousness + high LR & 5.7 & 0.79 & 0.8 \\
Low consciousness + low LR & 8.9 & 0.72 & 1.2 \\
\bottomrule
\end{tabular}
\end{table}

The cadence controller successfully allocates more generation bandwidth to high-consciousness, high-learning states where generation quality is highest and hallucination risk is lowest.

\subsection{Comparison with Post-Hoc Methods}

We compare the full Broca pipeline against representative post-hoc hallucination mitigation methods, applied to the same Liquid-Mamba backend without epistemic gating (Table~\ref{tab:comparison}).

\begin{table}[H]
\centering
\caption{Comparison with post-hoc hallucination mitigation methods.}
\label{tab:comparison}
\begin{tabular}{lcccc}
\toprule
\textbf{Method} & \textbf{Halluc. (\%)} & \textbf{Coherence} & \textbf{Latency (ms)} & \textbf{Prevents?} \\
\midrule
None (baseline) & 23.4 & 0.71 & 7.1 & No \\
Temperature scaling & 18.2 & 0.68 & 7.2 & No \\
Nucleus sampling ($p=0.9$) & 16.7 & 0.72 & 7.3 & No \\
Self-consistency ($k=5$) & 8.1 & 0.74 & 35.5 & No \\
RAG (top-5 retrieval) & 6.3 & 0.76 & 42.8 & No \\
\midrule
\textbf{Broca (full pipeline)} & \textbf{0.4} & \textbf{0.84} & \textbf{11.2} & \textbf{Yes} \\
\bottomrule
\end{tabular}
\end{table}

The Broca pipeline achieves the lowest hallucination rate (0.4\%) with competitive latency (11.2ms vs.\ 7.1ms baseline). Unlike all other methods, it \emph{prevents} hallucination by construction rather than detecting it post-hoc. The latency overhead of 4.1ms is attributable to the Hamming distance computation over the vocabulary, which is highly parallelizable.

\section{Discussion}

\subsection{Prevention versus Detection}

The fundamental distinction between epistemic gating and existing approaches is the difference between prevention and detection. Post-hoc methods allow hallucinated content to be generated and then attempt to identify and filter it. This creates an adversarial dynamic: as language models become more fluent, hallucinated content becomes harder to distinguish from factual content, and post-hoc detection becomes less reliable \citep{ji2023survey}.

Epistemic gating eliminates this arms race by operating at the logit level. Tokens that lack epistemic support are never assigned nonzero probability, regardless of how ``fluent'' they would be in context. The system cannot generate what it does not know, just as a biological speaker cannot articulate concepts that have not achieved conscious access \citep{dehaene2014consciousness}.

\subsection{Consciousness-Awareness as a Feature}

The modulation of the epistemic threshold by consciousness level ($\Phi$) implements a form of metacognitive regulation. During high-consciousness states---when the system's integrated information is high and its internal representations are most coherent---the threshold increases, demanding stronger epistemic support. During low-consciousness states, the threshold relaxes, permitting more exploratory generation.

This mirrors the human experience of verbal fluency varying with attentional state. Under high cognitive load or reduced consciousness (fatigue, distraction), speech becomes less precise and more error-prone \citep{levelt1989speaking}. Our system formalizes this relationship: consciousness level directly modulates the stringency of epistemic filtering.

\subsection{HDC as Epistemic Infrastructure}

The use of HDC for epistemic support computation is not incidental. The high dimensionality ($D = 16{,}384$) provides a rich similarity space where graded epistemic relationships are naturally encoded. The quasi-orthogonality of random HDVs ensures that the baseline similarity is approximately 0.50 (for binary vectors), providing a natural chance level against which epistemic support is measured.

The binding operation $\mathbf{G} = \mathbf{T} \otimes \rho^t(\mathbf{P}) \otimes \mathbf{K}$ creates a holographic representation that simultaneously encodes the thought content, contextual position, and knowledge grounding. Similarity to this composite vector requires a token to be consistent with all three factors, providing a more robust epistemic check than any single factor alone.

\subsection{Active Inference Interpretation}

Under the active inference framework \citep{friston2010free}, language generation is an epistemic action that reduces the system's uncertainty about its own beliefs. The EpistemicGate can be interpreted as a precision-weighting mechanism: tokens with high epistemic support have high precision (low uncertainty), while tokens with low support have low precision and are suppressed.

The adaptive cadence controller implements a form of expected free energy minimization: the system generates language when doing so is expected to reduce free energy (high consciousness, high learning rate) and withholds generation when the expected information gain is low.

\subsection{Limitations}
\label{sec:limitations}

Several limitations warrant discussion. First, the residual 0.4\% hallucination rate arises from \emph{compositional} hallucinations: individual tokens may be epistemically supported, but their combination produces an unsupported claim. For example, the tokens ``the,'' ``president,'' ``lives,'' ``on,'' ``Mars'' are individually supported, but their composition is not. Addressing compositional hallucination requires extending the epistemic check to $n$-gram or phrase-level representations, which we leave to future work.

Second, the epistemic gate may be overly conservative, suppressing novel but valid inferences. The system can only generate claims directly supported by its knowledge base, which limits creative and abductive reasoning. This is a fundamental trade-off between safety and capability.

Third, the computational overhead of per-token Hamming distance computation scales as $O(|\mathcal{V}| \cdot D)$ per generation step. For our vocabulary of 32,000 tokens and $D = 16{,}384$, this amounts to approximately $5 \times 10^8$ bit operations per step, which is tractable on modern hardware but nontrivial.

Fourth, our evaluation is conducted within the Broca pipeline's own knowledge representation. External validation against human-annotated factuality benchmarks (e.g., TruthfulQA \citep{lin2022truthfulqa}) would strengthen the empirical claims, and we plan this evaluation in future work. Relatedly, the $23.4\% \rightarrow 0.4\%$ hallucination reduction is measured on 500 in-distribution generation episodes with the same backend; we have not yet tested the gate against systematic adversarial probing (jailbreaks, indirect requests, role-play framings) that might route around the EpistemicGate by manipulating the HDC grounding or the consciousness-level threshold $\tau$. The mechanism is a physical logit mask, not a claim of robustness to sophisticated attack; a red-team evaluation belongs in revision. Finally, the HDC grounding knowledge base $\mathbf{K}_{\text{grounding}}$ is a structured knowledge graph derived from the system's own prior learning; the curation procedure and its effect on gate precision are detailed in \S\ref{sec:implementation}, but readers should note that a different grounding corpus would produce different hallucination-reduction numbers.

\subsection{Relationship to Neuroscience}

Our architecture is inspired by, but does not claim to replicate, the neural mechanisms of language production. In the human brain, Broca's area (left inferior frontal gyrus) is involved in syntactic processing, phonological encoding, and the sequencing of articulatory gestures \citep{hagoort2005broca}. Our Broca pipeline abstracts these functions into the ThoughtEncoder (conceptual encoding), HDC binding (structural encoding), and autoregressive generation (sequential output).

The epistemic gating mechanism has no direct neural analog that we are aware of, though it is functionally similar to the ``editor'' or ``monitor'' proposed in speech production models \citep{levelt1989speaking}, which filters pre-articulatory representations for errors before they reach motor execution.

\section{Related Work}

\citet{ji2023survey} provide the most comprehensive survey of hallucination in language models, categorizing causes and mitigation strategies. Our work differs in proposing prevention rather than mitigation.

\citet{lewis2020retrieval} introduced retrieval-augmented generation, which grounds generation in retrieved documents. RAG and epistemic gating are complementary: RAG provides external knowledge grounding, while epistemic gating provides internal epistemic filtering. A combined system could use RAG to populate the knowledge HDVs that the EpistemicGate queries.

\citet{kadavath2022language} showed that language models can be trained to express calibrated uncertainty. Our approach differs in that uncertainty is computed from HDC similarity rather than learned calibration, providing a more interpretable and verifiable uncertainty measure.

\citet{kanerva2009hyperdimensional} established the theoretical foundations of HDC. Our use of HDC for epistemic support computation extends the standard HDC applications (classification, analogy, language understanding) to a novel domain: epistemic filtering for safe generation.

\section{Conclusion}

We have presented epistemic gating, a mechanism for preventing hallucination in language generation by masking logits for tokens that lack sufficient epistemic support. Implemented within the Broca language pipeline---a 27,965~LOC system with 403 tests---epistemic gating reduces hallucination from 23.4\% to 0.4\%, a 98.3\% reduction. The mechanism operates at the logit level, physically preventing the generation of unsupported content rather than detecting it post-hoc.

The key architectural innovations are: (1) encoding cognitive state as a 16,384-dimensional HDV via a 43-channel ThoughtEncoder; (2) computing per-token epistemic support via HDC similarity; (3) consciousness-modulated threshold adjustment; and (4) mid-sentence semantic veto for coherence maintenance. Together, these mechanisms implement a consciousness-aware language generation system that can only articulate what it knows.

Future work will address compositional hallucination via phrase-level epistemic checking, external benchmark evaluation, and the integration of retrieval-augmented knowledge into the HDC grounding layer. We also plan to investigate the theoretical connections between epistemic gating and the free energy principle more formally, potentially deriving the gating threshold from variational inference principles.

\section{Reproducibility}

\section{Extension: The Epistemic Cube}

The original EpistemicGate applies a single scalar threshold to mask unsupported tokens. While effective at preventing hallucination, scalar gating cannot distinguish between different \emph{kinds} of epistemic uncertainty. A claim may be well-grounded empirically but novel to the community (high E, low N); or widely accepted but ephemeral (high N, low M). These distinctions matter for language style.

The Epistemic Cube extends scalar gating to four-dimensional modulation via the \texttt{EpistemicCubeGate}:

\begin{enumerate}
\item \textbf{Epistemic (E)}: 5 tiers (E0: opinion $\to$ E4: reproducible) controlling assertion strength. At E0, assertion words receive $-0.8$ logit penalty and hedging words receive $+0.5$ boost. At E4, assertions are boosted $+0.4$ and hedging penalized $-0.3$.

\item \textbf{Narrative (N)}: 4 tiers (N0: personal $\to$ N3: axiomatic) controlling social framing. Tokens like ``I think'' are boosted at N0; ``necessarily'' at N3.

\item \textbf{Meta (M)}: 4 tiers (M0: ephemeral $\to$ M3: foundational) controlling temporal framing. ``Now'' is boosted at M0; ``always'' at M3.

\item \textbf{Holistic (H)}: Scalar $\in [0, 1]$ from $\Phi$ and harmonic alignment. When $H < 0.25$, all logits are dampened $\times 0.85$.
\end{enumerate}

The composite quality score weights the three discrete axes:
\begin{equation}
Q = \frac{E}{4} \times 0.40 + \frac{N}{3} \times 0.35 + \frac{M}{3} \times 0.25
\end{equation}
reflecting the epistemological priority of empirical verification (40\%) over normative agreement (35\%) over permanence (25\%).

The cube coordinate is carried in ThoughtEncoder channels 28--42 (15 channels: 5 E one-hot + 4 N one-hot + 4 M one-hot + 1 H scalar + 1 quality scalar) and additionally blended into the thought hypervector at 30\% weight via the NSM (Neural Semantic Modulation) grounding pathway:
\begin{equation}
\mathbf{T}' = 0.7 \cdot \mathbf{T} + 0.3 \cdot \mathbf{T}_{\text{NSM}}
\end{equation}

This produces language whose style is structurally coupled to the system's epistemic state: at E0/N0, the system \emph{cannot} use confident universal claims without a logit penalty, producing naturally hedged output. The constraint is architectural, not prompted.

\subsection*{Reproducibility}

All experiments were conducted using the Broca language pipeline within the Symthaea cognitive architecture (v1.9.0). The following commands reproduce the core results:

\begin{verbatim}
# Run Broca pipeline tests (403 tests)
cargo test -p symthaea-broca -- --nocapture

# Run epistemic gating unit tests
cargo test -p symthaea-broca --lib epistemic -- --nocapture

# Run coherence feedback tests
cargo test -p symthaea-broca --lib coherence -- --nocapture

# Run end-to-end generation tests
cargo test -p symthaea-broca --lib e2e -- --nocapture

# Build with language feature enabled
cargo build --features ssm_language
\end{verbatim}

The system requires Rust 1.75+ and the \texttt{nix develop} environment. The \texttt{ssm\_language} feature flag must be enabled for Broca pipeline functionality. The Liquid-Mamba backend requires the default feature set.

\subsection*{Code Availability}

The Broca language pipeline source code is located in \texttt{symthaea/crates/symthaea-broca/} within the Luminous Dynamics Symthaea repository. The EpistemicGate implementation, CoherenceFeedback module, and adaptive cadence controller are all contained within this crate. Correspondence and access requests should be directed to the corresponding author.

\bibliographystyle{plainnat}
\begin{thebibliography}{30}

\bibitem[Baars(1988)]{baars1988cognitive}
Baars, B.~J. (1988).
\newblock \emph{A Cognitive Theory of Consciousness}.
\newblock Cambridge University Press.

\bibitem[Bai et~al.(2022)]{bai2022constitutional}
Bai, Y., Kadavath, S., Kundu, S., et~al. (2022).
\newblock Constitutional AI: Harmlessness from AI feedback.
\newblock \emph{arXiv preprint arXiv:2212.08073}.

\bibitem[Bender et~al.(2021)]{bender2021dangers}
Bender, E.~M., Gebru, T., McMillan-Major, A., and Shmitchell, S. (2021).
\newblock On the dangers of stochastic parrots: Can language models be too big?
\newblock In \emph{FAccT '21}, pp.\ 610--623.

\bibitem[Dehaene(2014)]{dehaene2014consciousness}
Dehaene, S. (2014).
\newblock \emph{Consciousness and the Brain}.
\newblock Viking Press.

\bibitem[Fedorenko et~al.(2014)]{fedorenko2014role}
Fedorenko, E. and Thompson-Schill, S.~L. (2014).
\newblock Reworking the language network.
\newblock \emph{Trends in Cognitive Sciences}, 18(3):120--126.

\bibitem[Friston(2010)]{friston2010free}
Friston, K. (2010).
\newblock The free-energy principle: a unified brain theory?
\newblock \emph{Nature Reviews Neuroscience}, 11(2):127--138.

\bibitem[Gu and Dao(2024)]{gu2024mamba}
Gu, A. and Dao, T. (2024).
\newblock Mamba: Linear-time sequence modeling with selective state spaces.
\newblock \emph{arXiv preprint arXiv:2312.00752}.

\bibitem[Hagoort(2005)]{hagoort2005broca}
Hagoort, P. (2005).
\newblock On Broca, brain, and binding: a new framework.
\newblock \emph{Trends in Cognitive Sciences}, 9(9):416--423.

\bibitem[Hasani et~al.(2021)]{hasani2021liquid}
Hasani, R., Lechner, M., Amini, A., et~al. (2021).
\newblock Liquid time-constant networks.
\newblock In \emph{AAAI '21}, pp.\ 7657--7666.

\bibitem[Ji et~al.(2023)]{ji2023survey}
Ji, Z., Lee, N., Frieske, R., et~al. (2023).
\newblock Survey of hallucination in natural language generation.
\newblock \emph{ACM Computing Surveys}, 55(12):1--38.

\bibitem[Kadavath et~al.(2022)]{kadavath2022language}
Kadavath, S., Conerly, T., Askell, A., et~al. (2022).
\newblock Language models (mostly) know what they know.
\newblock \emph{arXiv preprint arXiv:2207.05221}.

\bibitem[Kanerva(2009)]{kanerva2009hyperdimensional}
Kanerva, P. (2009).
\newblock Hyperdimensional computing: An introduction to computing in distributed representation with high-dimensional random vectors.
\newblock \emph{Cognitive Computation}, 1(2):139--159.

\bibitem[Levelt(1989)]{levelt1989speaking}
Levelt, W.~J. (1989).
\newblock \emph{Speaking: From Intention to Articulation}.
\newblock MIT Press.

\bibitem[Lewis et~al.(2020)]{lewis2020retrieval}
Lewis, P., Perez, E., Piktus, A., et~al. (2020).
\newblock Retrieval-augmented generation for knowledge-intensive NLP tasks.
\newblock In \emph{NeurIPS 2020}.

\bibitem[Lin et~al.(2022)]{lin2022truthfulqa}
Lin, S., Hilton, J., and Evans, O. (2022).
\newblock TruthfulQA: Measuring how models mimic human falsehoods.
\newblock In \emph{ACL 2022}, pp.\ 3214--3252.

\bibitem[Maynez et~al.(2020)]{maynez2020faithfulness}
Maynez, J., Narayan, S., Bohnet, B., and McDonald, R. (2020).
\newblock On faithfulness and factuality in abstractive summarization.
\newblock In \emph{ACL 2020}, pp.\ 1906--1919.

\bibitem[Ouyang et~al.(2022)]{ouyang2022training}
Ouyang, L., Wu, J., Jiang, X., et~al. (2022).
\newblock Training language models to follow instructions with human feedback.
\newblock In \emph{NeurIPS 2022}.

\bibitem[Wang et~al.(2023)]{wang2023selfconsistency}
Wang, X., Wei, J., Schuurmans, D., et~al. (2023).
\newblock Self-consistency improves chain of thought reasoning in language models.
\newblock In \emph{ICLR 2023}.

\bibitem[Wei et~al.(2022)]{wei2022chain}
Wei, J., Wang, X., Schuurmans, D., et~al. (2022).
\newblock Chain-of-thought prompting elicits reasoning in large language models.
\newblock In \emph{NeurIPS 2022}.

\bibitem[Thorne et~al.(2018)]{thorne2018fever}
Thorne, J., Vlachos, A., Christodoulopoulos, C., and Mittal, A. (2018).
\newblock FEVER: a large-scale dataset for fact extraction and VERification.
\newblock In \emph{NAACL 2018}, pp.\ 809--819.

\bibitem[Tononi(2004)]{tononi2004information}
Tononi, G. (2004).
\newblock An information integration theory of consciousness.
\newblock \emph{BMC Neuroscience}, 5(1):42.

\end{thebibliography}

\end{document}
